\documentclass[../main.tex]{subfiles}
\begin{document}
%==========================================TIÊU ĐỀ TẬP========================================
\begin{table}[h]
  \begin{adjustwidth}{-1cm}{}
    \begin{tabular}{>{\centering\arraybackslash}p{6cm}|>{\raggedright\arraybackslash}p{12.5cm}}
      \multirow{5}{6cm}
      % Cột bên trái: ÉPISODE và số tập
      {\\[-40pt]
        \flaregothic\fontsize{20pt}{20pt}\selectfont \centering
        \color{eptype}ÉPISODE\color{black}\\
        \burbank\fontsize{72pt}{72pt}\selectfont
        \color{epnum}122
        \\[18pt]
        % \flaregothic\fontsize{14pt}{14pt}\selectfont
        % \color{epattr}BIENVENUE, \uline{\color{Banpire} Mel} !
        % \flaregothic\fontsize{18pt}{18pt}\selectfont
        % \color{epattr}ÉPISODE PERDU
        \color{black}
      }
       &
      % Dòng thứ nhất: tên tập tiếng Nhật
      {\fotskip\fontsize{18pt}{18pt}\selectfont \ruby{凶}{きょう}\ruby{器}{き}\ruby{持}{も}たせてみた
      }                   \\[6pt]
      % Dòng thứ hai: tên tập tiếng Anh
       & {\cafeteria\fontsize{18pt}{18pt}\selectfont Sharp Objects}                                       \\[4pt]
      % Dòng thứ ba: tên tập tiếng Việt
       & {\vnmsans\fontsize{18pt}{18pt}\selectfont Chuyện về cái cây kim}                                 \\[8pt]
      % Dòng thứ tư: Nhân vật xuất hiện
       & {\sen{\fontsize{14pt}{14pt}\selectfont Track used: Knit Kingdom  (ニット\ruby{王}{おう}\ruby{国}{こく})}}
      \\[6pt]
      % Dòng cuối cùng: Thời gian diễn ra của tập (thường sẽ là ngày phát hành)
       & {\continuum\fontsize{16pt}{16pt}\selectfont Ngày phát hành: 12/09/2021}                          \\[6pt]
       & (dựa theo bản dịch của \href{https://www.facebook.com/mamisubteam}{MamiSubs - Mami Fansub})      \\[18pt]
    \end{tabular}
  \end{adjustwidth}
\end{table}
%=========================================TRUYỆN CHÍNH========================================
\color{black}

Gần đến trung tuần tháng Chín, bầu không khí ở văn phòng \hll\ trở nên mát mẻ và yên tĩnh đến lạ, khi mùa thu đã dần tỏa ra khắp phố.

Nhưng, ``sự yên tĩnh'' ở đây cũng chỉ\dots\ đúng ở một khía cạnh nào đó.

Ngay lúc này đây ở góc phòng, Rushia ngồi tỉ mẩn trên ghế bành, trên tay là một sợi chỉ nhỏ màu đỏ và một cây kim khâu lấp lánh ánh thép. Cô tập trung đến mức không chớp mắt, lưỡi thỉnh thoảng thò ra mép như đang lắp ráp bom hẹn giờ.

Đối diện cô, Kuon, Mel và Lamy đang nhìn cô như đang quan sát một sinh vật lạ đang ủ mưu thống trị thế giới với những vẻ mặt khó hiểu.

Cậu Tổng ngồi ở ghế bành đối diện, tay chống cằm lên tiếng: ``Rushia-chan, em đã như thế khoảng mười lăm phút rồi đấy\dots''

Nàng Dơi đứng ở phía sau, hơi trố mắt: ``Em ấy đang tập trung quá\dots\ giống như đang nhập hồn vậy à\dots''

Nàng Yêu tinh tuyết giật mình, quay sang Rushia, rồi về phía Kuon: ``Trước mặt chúng ta mà lại say sưa với cây kim như thế? Anh chị đã thử gọi chị ấy chưa?''

Kuon chỉ nhún vai đáp: ``Vừa nãy anh có gọi, nhưng em ấy chỉ bảo là: `Tập trung cho em làm, nanodesu' và thế là anh cứ để vậy thôi.''

Lamy gật gù, rồi ánh mắt ánh lên một tia ranh mãnh. Cô thì thầm: ``Phải đánh lạc hướng chị ấy thôi.''

``Ể?'' Mel giật mình, ngạc nhiên với ý tưởng có phần ngớ ngẩn này.

Cậu Tổng nghe vậy, liền lên tiếng căn ngăn: ``Anh\dots\ không nghĩ đó là ý hay đâu. Cứ kệ em ấy đi.''

``Nhỡ đâu chị ấy lại nghĩ đến mấy thứ gì đó kinh dị thì sao?'' nàng Yêu tinh tuyết hỏi lại, vẻ mặt nửa lo lắng, nửa ranh ma.

Cậu nhìn sang Rushia vẫn đang lẩm bẩm gì đó với kim và chỉ, rồi đáp lại cẩn trọng: ``\dots Xét về tính cách yandere của Ru-chan, chắc em nói cũng đúng đấy. Nhưng mà anh nghĩ hai đứa chỉ cần gọi em ấy thôi\dots''

Dù đã cảnh báo, Mel và Lamy vẫn quyết định hành động. Kuon thở dài, giơ hai tay đầu hàng: ``Em ấy mà nổi cáu là anh không chịu trách nhiệm đâu đấy.''

\ct

Lamy bước ra phía sau Rushia, lấy hơi thật sâu rồi hú lên một tiếng ``ma ám'' giả trân: ``Uoooaahh\dots\~{}\~{}\~{}''

Rushia không hề phản ứng. Cô vẫn nhìn chằm chằm vào đầu kim, không để ý gì tới xung quanh.

``\hl{Em ấy trông có vẻ tập trung cao độ quá\dots}'' cậu lẩm bẩm, tiếp tục ngồi xem nàng Pháp sư.

Nhưng hai cô gái còn lại vẫn không hề cam chịu. Mel liền nhào ngược người lên trần nhà, treo ngược bằng chân lên lan can tầng lửng, hai tay dang rộng như dơi thật sự. ``Rushiaaa\~{} Nhìn lên đây nèeee\~{}\~{}''

\img{10-2a}

Vẫn không có phản ứng nào. Cả ba người bắt đầu đổ mồ hôi lạnh.

Nàng Yêu tinh tuyết than thở: ``Không được rồi\dots\ Chị ấy tập trung quá!''

Kuon ngán ngẩm nói: ``Thôi, cứ ngơ em ấy đi. Tí nữa xong rồi gọi em ấy cũng được.''

Nghĩ rằng nếu tiếp cận theo cách dọa dẫm thì không có ý nghĩa gì, Lamy bất chợt nảy ra một ý tưởng khác. Cô nhíu mày: ``Kuon-senpai, chị ấy sợ cái gì nhất?''

Câu hỏi có phần bất chợt đó làm cậu khẽ khựng lại, nhưng rồi cũng xoa cằm suy nghĩ và đáp: ``Anh nghĩ là một Cô đồng thì không sợ ma quỷ hay gì đâu\dots\ Nhưng mà con gái thường hay sợ  sâu bọ hoặc là ếch nhái\dots\ chắc thế?''

Lamy vỗ tay cái bốp như vừa nghĩ ra ý tưởng: ``Vậy để em thử.''

Ngay lập tức, cô cúi sát tai Rushia, bất ngờ hét lên: ``\textbf{Sau lưng chị có ếch kìa, Rushia-senpai!!}''

``\textbf{ÁÁÁÁÁÁÁÁÁÁÁÁÁÁÁÁÁÁÁÁ!!!!}''

Lập tức, nàng Pháp sư trợn ngược mắt lên, mặt tái mét, bật dậy và vô thức ném thứ đang cầm trên đi như ninja ném phi tiêu. Kuon và Mel cùng bật dậy theo bản năng.

``Vãi chưởng\dots!'' cậu Tổng giật mình, bật ngửa ra sau, còn nàng Dơi thì đổ mồ hôi: ``C-Có hiệu quả thật kìa\dots''

Nhưng dường như cậu vừa nhận ra thứ cô ném. Cậu đập tay vào trán, thốt ra những lời mà cậu đáng ra nên nói từ ba giây trước: ``\dots Này! Đừng có ném cái kim đi chứ!!''

Cùng lúc đó, một tiếng *CẠCH* vang lên trước cửa văn phòng.

Cánh cửa bật mở. Kanata hớn hở bước vào, tay giơ cao túi bánh, miệng chuẩn bị hét lên một câu chào tươi rói: ``Chào buổi sá-''

*PHẬP!*

\begin{figure}[H]
  \centering
  \begin{subfigure}[t]{0.45\textwidth}
    \centering
    \includegraphics[width=8cm]{../../../../Image/Book?/10-2b1.png}
    \caption{Rushia phóng cây kim đi vì hoảng sợ\dots\dots}
  \end{subfigure}
  \quad
  \begin{subfigure}[t]{0.45\textwidth}
    \centering
    \includegraphics[width=8cm]{../../../../Image/Book?/10-2b2.png}
    \caption{\dots trúng vào cánh cửa văn phòng, ngay sát đầu của Kanata.}
  \end{subfigure}
\end{figure}

Cây kim sắc lẹm cắm thẳng vào cánh cửa ngay sát mang tai cô, rung lên bần bật nhẹ, tạo thành một luồng sóng âm vì lực quá mạnh. Nụ cười trên mặt Thiên sứ tắt lịm, đầu giật giật nhìn về cây kim, chân khuỵu xuống.

``L-Là em mà\dots\ Amane Kanata đây\dots'' giọng cô run lên, mồ hôi túa ra như tắm. Rồi, cô gục xuống ngay sau đó.

Kuon lao ra đỡ kịp lúc, tựa cô vào cửa: ``Em có sao không vậy, Kana-chan!?''

Rushia cuống cuồng chạy đến rối rít: ``C-Chị xin lỗi, Kanatan!! Chị không cố ý!!''

Mel nhìn Kanata đã bay hồn bạt vía, tỏ vẻ thương cảm cho nàng hậu bối suýt trở thành nạn nhân kia: ``Trời ạ\dots\ Mặt em ấy như thể hồn lìa khỏi xác luôn\dots''

Cô quay sang nàng Thế hệ thứ 3, tò mò: ``Cơ mà\dots\ em không biết Rushia-chan lại sợ ếch đến thế\dots''

Cậu Tổng thở dài, nhún vai đáp: ``Ờ thì\dots\ con gái sợ mấy con sâu bọ hay lưỡng cư là chuyện bình thường thôi mà\dots\ Cái này thì bây giờ anh mới biết đấy.''

\ct

Mel vừa dìu Kanata vào phòng vừa nhẹ nhàng xoa đầu cô như vỗ về em bé. Thiên sứ tóc bạch kim vẫn còn tái mặt, hai tay ôm đầu, mắt thì mơ màng như vẫn chưa rút được hồn về.

Kuon khoanh tay lại, thở dài: ``Suýt tí nữa là em hại người ta rồi đấy\dots''

Rushia ngồi phịch xuống sofa, vừa vén lại lớp váy bị xô lệch, vừa lầm bầm phản pháo: ``Em xin lỗi rồi mà\dots\ C-cơ mà\dots'' Cô liếc mọi người một cái đầy bực bội, tiếp tục càu nhàu: ``Em đang tập trung đó! Mọi người trật tự đi.''

Nói rồi cô đứng dậy, đi về phía cửa, rút cây kim đang cắm vào cánh gỗ như rút dao khỏi bao, rồi nghiêm túc quay lại tiếp tục ``sự nghiệp'' xâu chỉ của mình.

Lamy nghiêng đầu, hỏi với vẻ tò mò thật sự: ``Chị đang may đồ hả?''

Nàng Pháp sư cau mày, mắt vẫn tập trung vào kim chỉ: ``May mủng cái gì. Xâu chỉ luồn kim thôi mà.''

Cả căn phòng bỗng chốc im bặt sau câu trả lời có phần cụt lủn từ cô ấy.

``\dots Thế thôi á?'' cậu con trai lên tiếng, nhíu mày nhẹ.

Nàng Dơi vẫn đang vuốt đầu Kanata, gặng hỏi: ``Nhưng mà em làm thế để làm gì?''

Rushia ngẩng đầu lên, ánh mắt sáng long lanh như một cô nàng mơ mộng đang nói triết lý sống: ``Để luyện thêm khả năng tập trung cao độ và khéo léo! Và còn-'' cô đứng dậy, giơ gtay tạo dáng như người mẫu cho quảng cáo, ``-trở thành một cô gái toàn năng!''

\img{10-2c}

Nàng Thiên sứ - giờ đã hoàn hồn dù vẫn đang tựa vào Mel - nhăn mặt: ``Làm thế cũng không toàn năng thêm được nữa đâu chị\dots''

Cậu Tổng thì bật cười nhẹ, lắc đầu với lời giải thích bất hợp lý này: ``Anh thấy xâu chỉ chả liên quan gì đến việc toàn năng cả. Cơ mà nếu em bảo là để `rèn luyện sự khéo léo với tập trung cao độ' thì còn chấp nhận được.''

Nàng Yêu tinh tuyết lúc này chống hông, phồng má: ``Gì chứ, cái này em nhắm mắt cũng làm được!''

``Hả!?'' Rushia quay đầu lại, trợn mắt nghi ngờ. Kuon thì nheo mắt nhìn cô, bĩu môi: ``Kinh thế! Giỏi thì làm thử đi xem nào!''

Lamy chẳng chần chừ. Cô rút lấy cây kim, cầm sợi chỉ, nhắm tịt mắt, rồi hô vang: ``Xem em đây!''

Cô hướng tay cầm cây kim về phía tay cầm chỉ, và bởi vì cô không hề hay biết rằng cô đang nhầm tay, nên\dots

*PHẬP*

\dots việc bị cây kim đâm thẳng vào là điều không thể tránh khỏi.

\begin{figure}[H]
  \centering
  \includegraphics[width=12.5cm]{../../../../Image/Book?/10-2d.png}
  \caption{Lamy bị ngược tay: thay vì đưa tay cầm chỉ xâu lỗ kim, cô lại\dots\ hướng tay đang cầm kim vào sợi chỉ, kết quả như ta thấy.}
\end{figure}

``\textbf{Ái daaaaa!!}'' Lamy la lớn, rút ngón tay lại như bị bỏng. Khuôn mặt cô xanh mướt, trực trào nước mắt vì đau.

Kuon xoa cằm, mỉa mai bằng giọng chưng hửng: ``Làm còn lộn tay nữa. Đấy, đã ngốc lại còn thích thể hiện ta đây cơ.''

Rushia lớn giọng quát: ``\textbf{Bộ não em có vấn đề à!?}''

Kanata thì thở dài, lắc đầu ngán ngẩm: ``Thật. Em mở mắt còn chưa làm được đây này.''

Lamy rơm rớm nước mắt, đưa ngón tay lên thổi phù phù: ``Chảy máu rồi\dots\ huhu\dots''

Bất thình lình, Mel - ngồi ngay sau Lamy - run bắn cả người, ôm đầu lùi lại hai bước, mặt xám xịt như bầu trời Hà Nội vào mùa hè.

Nàng Thiên sứ tròn mắt phản ứng: ``Chẳng phải chị là ma cà rồng sao!? Sao lại sợ máu chứ!?''

\img{10-2e}

Trong lúc hỗn loạn, Rushia lại tiếp tục xâu chỉ, nhưng sau vài giây chật vật vẫn chưa làm được. Cô thở dài, vẻ mặt đượm buồn: ``Nhưng mà công nhận\dots\ Cái này khó quá à\dots''

Kuon nhìn nàng Pháp sư cực nhọn, hết trật chỗ này tới chỗ khác, cũng phải gật gù: ``Quan trọng ở đây là phải kiên nhẫn. Em cứ xỏ bừa bãi thế không trúng đâu.''

Nàng Dơi lúc này bõng chốc vỗ tay nảy ra một ý tưởng khác: ``Em phải đặt mình vào vị trí của sợi chỉ!''

Nàng Yêu tinh tóc trắng ở sau, với ngón tay được băng bó, giơ hai nắm đấm lên phụ họa: ``Chính\~{} xác\~{}!''

``Cụ thể là thế nào?'' / ``Gì chứ?'' cả cậu Tổng và nàng Pháp sư đều cùng ngơ ngác nhìn cả hai.

Không để bất cứ ai nói gì thêm, Mel lôi ra một cây kim siêu to khổng lồ dựng thẳng lên sàn nhà (không biết bằng cách nào), rồi hô lớn: ``Hãy luồn bản thân qua cây kim này!''

\img{10-2f}

Cậu Tổng đờ người ra, nhìn cây kim khổng lồ một lần nữa, rồi nhướng mày: ``\dots Sao anh thấy có gì đó sai sai vậy?''

Nàng Pháp sư cũng tỏ hoang mang không kém: ``Chị muốn em trở thành sợi chỉ à?''

Lamy và Kanata đồng loạt gật đầu, đưa tay ra mời.

\begin{dialogue}
  \speak{Lamy} Bọn em cũng sẽ giúp chị.
  \speak{Kanata} Ừ, cùng hợp sức nào!
\end{dialogue}

Rushia chớp mắt vài lần. Trong khoảnh khắc, cô cảm thấy\dots\ được yêu thương. Đôi mắt cô rưng rưng, cảm tưởng như mọi người hiểu được những khó khăn (tí tị) mà cô đang trải qua.

Chỉ có Kuon vẫn đứng đó, tay khoanh lại, mặt lạnh như tiền. Nhưng cậu cũng không nói gì. Cậu biết rằng, có ngăn cũng chẳng được, nên chỉ im lặng quan sát.

Hai cô gái mỗi bên đỡ một bên nách Rushia, rồi nâng cả người cô lên theo chiều dọc như một sợi chỉ sống. Kuon nhìn cảnh tượng kỳ lạ ấy, mặt đơ ra: ``\textit{Sao mình thấy cái này quen quen\dots}''

\img{10-2g}

``\textbf{Lên nào, mấy đứa!}'' nàng Cô đồng la lớn, cảm giác như một chiến binh Sparta. Ngay sau đó, ba cô nàng hét lên đồng thanh như sắp vào trong lòng chiến trườn của giặc.

``Ấy từ từ, ở dưới chân có\dots'' cậu Tổng hốt hoảng lên tiếng cảnh báo, nhưng\dots

*VÉO!*

\dots không kịp. Hai cô gái trượt phải vỏ chuối ai đó ăn xong vứt lại, và cả ba lao đi như tên lửa. Chỉ vài tích tắc sau:

*XOẸT! XOẸT! XOẸT!*

\gif{10-2h}{1}{52}

Lamy đáp trọn vào đáy của khe kim, theo sau là Rushia đáp trọn xuống đỉnh đầu của người trước, rồi Kanata chui vừa khít cuối cùng như một chiếc xiên thịt hoàn hảo.

Kuon há hốc mồm, chỉ tay: ``Xuyên qua thật luôn kìa! Lại qua đầu cả ba đứa nữa chứ!''

Mel đứng đơ người ra, rồi đập nắm đấm vào tay kia, tươi cười: ``Ba nàng\dots\ xiên dango!''

Ba giây trôi qua.

Không ai nói gì. Cả ba người đang ``xuyên kim'' vẫn chưa hiểu chuyện gì xảy ra. Rushia nhăn mặt: ``Thật là không cần thiết.''

Cậu Tổng gật đầu lia lịa: ``Ừ. Đồng ý. Thừa thãi vãi chưởng.''

Tới cuối cùng, nó vẫn không giải quyết được vấn đề mà nàng Pháp sư đang gặp phải vốn rất cỏn con: làm thế nào để xâu được vào lỗ kim.
%==========================================TRUYỆN PHỤ=========================================
\omk

Bộ ``ba cô nàng xiên dango'' vẫn đang chui vắt vẻo qua cái khe kim khổng lồ. Rushia thì có vẻ đang cựa quậy khó chịu, trong khi Lamy và Kanata lại thoải mái dựa người vào nhau, cười đùa vui vẻ như vừa hoàn thành một nghi thức trưởng thành nào đó.

Mel, khoanh tay đầy đạo mạo, đứng trước họ như một giáo viên đang giảng bài, nhẹ giọng: ``Hãy chắc rằng nỗ lực của mình phải đặt đúng chỗ. Vậy em đã học được gì chưa, Rushia-chan?''

Hẳn nhiên với sự khó chịu còn bứt rứt trong mình, cô gằn giọng đáp: ``Nó thực sự vẫn chẳng giúp ích được gì cho em cả-nanodesu.''

Nàng Yêu tinh tuyết lắc người nhẹ, vẫn đang mắc trong khe kim: ``Em lại thấy vui mà!''

Nàng Thiên sứ giơ hai tay như học sinh giơ bảng điểm tốt: ``Đúng đó! Ta nên làm lại đi!''

Kuon, vẫn khoanh tay đứng từ phía xa, lắc đầu, giọng đều đều pha chút mỉa mai: ``Anh thấy cách của em cũng không ổn đâu, Mel-chan ạ. Với cả\dots\ nó trật khỏi vấn đề cốt lõi rồi.''

Rushia giơ tay chỉ thẳng lên trời, thẳng thắn như một đại biểu đang trả lời phần chất vấn: ``Đúng đó! Cái em cần là xâu chỉ vào kim mà em vẫn chưa làm được này!''

Cậu Tổng nhìn cô vài giây, rồi ánh mắt bỗng đổi tông. Cậu nghiêng đầu, tay chống cằm như vừa nảy ra ý tưởng: ``Nhưng mà cách làm của Mel-chan lại có dụng ý đấy. Cứ hiểu theo cách khác là được.''

Nàng Pháp sư nhìn cậu, mặt trơ ra: ``Ể?''

Cậu không giải thích gì thêm, chỉ vẫy tay ``Đưa cho anh cái kim với sợi chỉ.''

Mel ngoan ngoãn đưa đồ qua tay cậu. Kuon cầm lấy, rồi chỉ nhẹ vào tay của cô nàng tóc xanh: ``Lúc anh xem em se chỉ, anh thấy em run tay nãy giờ. Thế nên toàn trật lất cũng phải.''

Rushia gãi gãi má, thẹn thùng: ``Thì\dots\ Đây là lần đầu em làm mà-nanodesu\dots''

Cậu gật đầu, tiếp tục giảng giải: ``Và vì run tay nên sợi chỉ cứ bị trượt. Anh nghĩ giải pháp này khả dĩ hơn\dots''

Cậu đảo mắt, tìm quanh văn phòng, rồi lôi ra một chiếc hộp bìa cứng nhỏ, loại dùng để đựng giấy in. Đặt nó xuống trước mặt, cậu tiếp tục: ``Kiếm cho anh một cái gì đó vững vững. Xong rồi\dots''

Với cây kim trên tay, cậu cắm thẳng xuống chiếc hộp, vừa đủ để nó đứng vững. ``Rồi, cứ để yên như thế.''

``Có cần thiết phải\dots\ cắm thẳng không?''

``Nghiêng cũng được, nhưng mà đừng nghiêng quá. Quan trọng là nó phải đứng yên. Rồi em nhắm một mắt lại, hít sâu, rồi\dots''

Cậu thao tác tay nhanh nhẹn, xâu vào lỗ khoảng ba lần trước khi sợi chỉ luồn vào.

*XOẸT*

``Đấy, đơn giản thế thôi.''

Chỉ trong vài giây, cậu đã xâu chỉ vào kim một cách gọn gàng, nhanh gấp mười lần so với những nỗ lực đầy kịch tính ban nãy.

Cả ba cô nàng dango và Mel đều há hốc mồm, đôi mắt sáng lấp lánh.

Rushia sững sờ ngạc nhiên khi cậu có xâu chỉ một cách dễ dàng: ``Thế thôi sao?''

Kuon gật đầu bình thản: ``Ừ. Chỉ thế. Nếu run thì cứ dùng mẹo này.''

``Thế mà em cứ tưởng-nanodesu\dots''

Lamy chống cằm nhìn cậu, ánh mắt ngưỡng mộ: ``Anh có nhiều mẹo vặt hay nhỉ\dots''

Kanata khoanh tay, gật gù: ``Em nghe đồn là anh ấy có biết vài mẹo vặt để làm ảo thuật từ Nene-chan\footnote{Xem lại {\flaregothic ÉPISODE 117}, quyển số Chín.} đấy!''

Cậu Tổng cười khẽ, khiêm tốn đáp: ``Vấn đề nằm ở sự kiên nhẫn cơ. Cứ nỗ lực nhiều vào là được. Về vấn đề này thì\dots''

Cậu quay về phía Mel, nghiêm túc hơn: ``Mel-chan nói cũng phần đúng đấy. Nếu biết đặt đúng chỗ thì sẽ thành công.''

Nàng Yêu tinh tuyết nhìn lại cây kim cắm trên hộp đã luồn chỉ thành công, chớp mắt: ``Nhưng hàm ý của anh cũng có nghĩa là\dots\ tìm đủ mọi cách để đạt được mà nhỉ?''

Nàng Dơi nháy mắt: ``Chị cũng nghĩ thế.''

Kuon chỉ đáp lại bằng một cái nhún vai: ``\dots Cũng tùy vào từng hoàn cảnh thôi.''

Cả nhóm đồng loạt cười phá lên. Chỉ có Rushia, lúc này đã bắt đầu mỏi, thều thào lên tiếng, mặt đỏ bừng: ``A-A-nồ\dots\ Mọi người có thể\dots\ ra được không? Nằm ở tư thế này em thấy ngột ngạt quá-nanodesu\dots''

Mel và Kuon phì cười. Cậu con trai xua tay: ``À, ừ, quên mất. Xin lỗi em nha, Rushia-chan!''

Cả hai nhanh chóng gỡ ba ``cô nàng dango'' khỏi khe kim khổng lồ, đặt họ xuống mặt đất an toàn. Rushia lúc này đã hơi đỏ mặt vì xấu hổ, nhưng trên môi lại khẽ mỉm cười.

Ngày hôm đó, không chỉ Cô đồng tóc xanh học được mẹo se chỉ, mà còn (vô tình) thu nhặt thêm một bài học cuộc sống vốn\dots\ chẳng liên quan gì cả.

Nhưng dẫu sao, cô vẫn cảm tạ hai người kia một cách thật lòng.

Kuon nhìn Rushia cùng ba cô gái tiếp tục tập xâu kim từ phía xa, khoanh tay thở dài với một nụ cười trên môi: ``\textit{Cơ mà\dots\ để trở thành một người đa-zi-năng, thì mấy đứa vẫn phải học nhiều đấy.}''

\end{document}